\section{XMSS}
SHRINCS employs two distinct Merkle tree constructions: a standard balanced XMSS tree for the stateless hypertree, and an Unbalanced XMSS (UXMSS) tree optimized for the stateful component (specified in Section~8).

The stateless signature mode uses a SPHINCS+-style hypertree consisting of \texttt{d} layers, where each layer contains standard balanced XMSS trees. Each XMSS tree authenticates \texttt{2\string^h'} WOTS+C key pairs, where \texttt{h' = hsl / d} is the height of trees within each layer.

\subsection{XMSS Tree Structure}
An XMSS tree of height \texttt{h'} is a complete binary Merkle tree where:
\begin{itemize}
    \item Leaves: The \texttt{2\string^h'} leaves are WOTS+C public keys, computed via \texttt{wots\_PKgen} with mode \texttt{SL}.
    \item Internal nodes: Each internal node is the tweakable hash of its two children, using address type \texttt{SL\_TREE}.
    \item Root: The tree root serves as either the layer's contribution to the hypertree or (for the top layer) as \texttt{PK.sl}.
\end{itemize}

Layer indexing within the stateless hypertree is the following:
\begin{itemize}
    \item Layer 0: Bottom layer (is used to sign the FORS+C public keys)
    \item Layer d-1: Top layer (root is \texttt{PK.sl})
\end{itemize}

\subsection{XMSS TreeHash}
The \texttt{xmss\_treehash} function computes subtree roots within an XMSS tree. It takes a target height and a leftmost leaf index, then recursively builds the subtree by computing leaves (WOTS+C public keys) at height 0 and hashing pairs of children to form internal nodes. This function is the core building block used both for computing the full tree root (by calling it with \texttt{target\_height = h'} and \texttt{start\_idx = 0}) and for generating individual sibling nodes needed in authentication paths (with arbitrary \texttt{start\_idx}).

\begin{verbatim}
    Algorithm: xmss_treehash(SK.seed, PK.seed, ADRS, target_height, start_idx)
    Input:
        SK.seed: secret seed
        PK.seed: public seed
        ADRS: tweak
        target_height: height of output node (0 = leaf level)
        start_idx: leftmost leaf index in the subtree
    Output:
        node: hash value

    1. if target_height == 0:
        a. pk ← wots_PKgen(SK.seed, PK.seed, ADRS, start_idx, SL)
        b. return pk
    2. left ← xmss_treehash(SK.seed, PK.seed, ADRS, target_height - 1, start_idx)
    3. right ← xmss_treehash(SK.seed, PK.seed, ADRS, target_height - 1, start_idx + 
                                                            + 2^(target_height - 1))
    4. setTypeAndClear(ADRS, SL_TREE)
    5. setTreeHeight(ADRS, target_height)
    6. setTreeIndex(ADRS, start_idx >> target_height)
    7. return Tw(PK.seed, ADRS, left || right)
\end{verbatim}

\paragraph{Note:} \texttt{start\_idx} need not be zero. When computing a full tree root, it is called with \texttt{start\_idx = 0}. When computing a sibling node for an authentication path, \texttt{start\_idx} is the leftmost leaf index of that sibling's subtree.

\subsection{XMSS Root Computation}
The \texttt{xmss\_root} function computes the root of an entire XMSS tree by invoking \texttt{xmss\_treehash} with the full tree height and starting from leaf index 0. The resulting root hash represents the public key for this XMSS tree instance, which either becomes \texttt{PK.sl} (for the top hypertree layer) or is signed by the layer above (for lower layers).
\begin{verbatim}
    Algorithm: xmss_root(SK.seed, PK.seed, ADRS, h_prime)
    Input:
        SK.seed: secret seed
        PK.seed: public seed
        ADRS: tweak
        h_prime: tree height
    Output:
        root: tree root (n bytes)

    1. return xmss_treehash(SK.seed, PK.seed, ADRS, h_prime, 0)
\end{verbatim}

\subsection{XMSS Authentication Path Generation}
For a leaf at index \texttt{idx}, the authentication path consists of the sibling nodes at each level from the leaf up to (but not including) the root. The authentication path contains \texttt{h\_prime} nodes, each of size \texttt{n} bytes, allowing a verifier to reconstruct the path from the leaf to the root.

At height \texttt{h}, the sibling subtree's leftmost leaf index is determined by flipping the h-th bit of \texttt{idx} (using XOR) and clearing all lower bits. The sibling node is then computed by calling \texttt{xmss\_treehash} on the subtree rooted at that position.
\begin{verbatim}
    Algorithm: xmss_auth_path(SK.seed, PK.seed, ADRS, h_prime, idx)
    Input:
        SK.seed, PK.seed: seeds
        ADRS: tweak
        h_prime: tree height
        idx: leaf index being authenticated
    Output:
    auth: authentication path (h_prime * n bytes)

    1. auth ← []
    2. for h from 0 to h_prime - 1:
        a. sibling_start ← (idx XOR (1 << h)) AND (0xFFFFFFFF << h)
        b. auth ← auth || xmss_treehash(SK.seed, PK.seed, ADRS, h, sibling_start)
    3. return auth
\end{verbatim}
\paragraph{Note:} The expression \texttt{(idx XOR (1 << h)) AND (0xFFFFFFFF << h)} flips the \texttt{h}-th bit of \texttt{idx} to identify the sibling, then clears the lower \texttt{h} bits to obtain the leftmost leaf index of the sibling's subtree of height \texttt{h}.

\subsection{XMSS Public Key Recovery from Signature}
Given a WOTS+C signature and authentication path, this algorithm recovers the XMSS tree root. First, it recovers the WOTS+C public key (the leaf value) from the WOTS+C signature. Then, starting from the leaf, it iteratively combines the current node with the corresponding authentication path node to compute the parent, continuing until reaching the root. At each height \texttt{h}, the \texttt{h}-th bit of idx determines whether the current node is the left child \texttt{(bit = 0)} or right child \texttt{(bit = 1)}, which dictates the ordering when hashing the pair.

The message \texttt{m} being verified is either a FORS+C public key (at layer 0) or a lower-layer tree root (at layers \texttt{1...d-1}). In both cases, m is already a hash, so \texttt{wots\_pkFromSig} is called with \texttt{is\_internal = true} (phase 1 of message hashing is skipped; see Section~6.7).

\begin{verbatim}
    Algorithm: xmss_pkFromSig(wots_sig, auth, m, PK.seed, PK.root, ADRS, h_prime, idx)
    Input:
        wots_sig: WOTS+C signature
        auth: authentication path (h_prime * n bytes)
        m: signed message (FORS+C public key or lower-layer tree root)
        PK.seed: public seed
        PK.root: root public key
        ADRS: tweak
        h_prime: tree height
        idx: leaf index
    Output:
        root: computed tree root (n bytes), or "false" if invalid

    1. pk ← wots_pkFromSig(wots_sig, m, PK.seed, PK.root, ADRS, idx, SL, true)
    2. if pk == false:
        return false
    3. node ← pk
    4. for h from 0 to h_prime - 1:
        a. setTypeAndClear(ADRS, SL_TREE)
        b. setTreeHeight(ADRS, h + 1)
        c. setTreeIndex(ADRS, idx >> 1)
        d. auth_node ← auth[h * n : (h + 1) * n]
        e. if (idx AND 1) == 0:
                node ← Tw(PK.seed, ADRS, node || auth_node)
            else:
                node ← Tw(PK.seed, ADRS, auth_node || node)
        f. idx ← idx >> 1
    5. return node
\end{verbatim}

\subsection{XMSS Signature Generation (Stateless Layer)}
To sign a message within the stateless hypertree, the XMSS signature generation produces a WOTS+C signature on the message followed by the authentication path for the corresponding leaf. The message being signed is typically either a FORS+C public key (at layer 0) or the root of a lower-layer XMSS tree (at layers 1 through \texttt{d-1}). In both cases the message is already a hash value, so \texttt{wots\_sign} is called with \texttt{is\_internal = true}.

The leaf index \texttt{idx} determines which WOTS+C key pair is used; this index is determined from the message digest during stateless signing (see Section~10.5).

\begin{verbatim}
    Algorithm: xmss_sign(m, SK.seed, SK.prf, PK.seed, PK.root, ADRS, h_prime, idx)
    Input:
        m: message to sign (lower-layer root or FORS public key)
        SK.seed: secret seed
        SK.prf: PRF key
        PK.seed: public seed
        PK.root: root public key
        ADRS: tweak
        h_prime: tree height
        idx: leaf index to use

    Output:
        sig: XMSS signature (WOTS+C signature || authentication path)

    1. wots_sig ← wots_sign(m, SK.seed, SK.prf, PK.seed, PK.root, ADRS, idx, SL, true)
    2. auth ← xmss_auth_path(SK.seed, PK.seed, ADRS, h_prime, idx)
    3. return wots_sig || auth
\end{verbatim}

\subsection{Contribution to Signature Size and Verification Complexity}
Each XMSS signature consists of a WOTS+C signature plus an authentication path of \texttt{h'} nodes: \texttt{R\_SIZE + c + (l + h') * n} bytes (484 bytes for SHRINCS-B and 1252 bytes for SHRINCS-L).

The full stateless hypertree consists of \texttt{d = 2} layers, giving:
\begin{itemize}
    \item 968 bytes and 4106 hashes for SHRINCS-B
    \item 2504 bytes and 130 hashes for SHRINCS-L
\end{itemize}

\paragraph{Note:} This excludes the FORS+C contribution (see Section~9). The total stateless signature size and verification cost are the sum of FORS+C and hypertree components, as presented in Section~3.2.2.